\documentclass[11pt,a4paper]{article}

\usepackage[T1]{fontenc}
\usepackage[utf8]{inputenc}
\usepackage[italian]{babel}
\usepackage{lmodern}
\usepackage{microtype}
\usepackage[a4paper,margin=2.5cm,headheight=15pt]{geometry}
\usepackage{amsmath,amssymb}
\usepackage{booktabs,longtable,tabularx,array}
\usepackage{enumitem}
\usepackage{xcolor}
\usepackage[most]{tcolorbox}
\usepackage{fancyhdr}
\usepackage{hyperref}
\usepackage{xurl}
\usepackage{listings}
\usepackage{tikz}
\usetikzlibrary{arrows.meta,positioning,calc}
\usepackage{titlesec}

\definecolor{DDTADark}{HTML}{1F2933}
\definecolor{DDTAGray}{HTML}{F2F4F7}
\definecolor{DDTABorder}{HTML}{8A97A3}
\definecolor{DDTABlue}{HTML}{2A5887}
\definecolor{DDTAGreen}{HTML}{EAF6EF}
\definecolor{DDTAGreenBorder}{HTML}{3F7D56}
\definecolor{DDTAWarn}{HTML}{FFF4D6}
\definecolor{DDTAWarnBorder}{HTML}{B7791F}
\definecolor{DDTAInfo}{HTML}{EAF3F8}
\definecolor{DDTAInfoBorder}{HTML}{3C718C}

\hypersetup{
  colorlinks=true,
  linkcolor=DDTADark,
  citecolor=DDTADark,
  urlcolor=DDTABlue,
  pdftitle={DDTA S1.5 - Requirement abstraction closure and provenance constraints},
  pdfauthor={Documentation-Driven Threat Analysis research workspace}
}

\pagestyle{fancy}
\fancyhf{}
\fancyhead[L]{\small Documentation-Driven Threat Analysis}
\fancyhead[R]{\small S1.5 closure revision}
\fancyfoot[C]{\thepage}
\renewcommand{\headrulewidth}{0.4pt}

\titleformat{\section}{\Large\bfseries\color{DDTADark}}{\thesection}{0.7em}{}
\titleformat{\subsection}{\large\bfseries\color{DDTADark}}{\thesubsection}{0.7em}{}
\setlist[itemize]{leftmargin=*,itemsep=0.28em,topsep=0.35em}
\setlist[enumerate]{leftmargin=*,itemsep=0.28em,topsep=0.35em}
\setlength{\emergencystretch}{3em}

\newtcolorbox{statusbox}[1]{
  enhanced,breakable,colback=DDTAInfo,colframe=DDTAInfoBorder,
  boxrule=0.8pt,arc=1.5mm,left=3mm,right=3mm,top=2.5mm,bottom=2.5mm,
  title={#1},fonttitle=\bfseries\sffamily,fontupper=\small\sffamily
}
\newtcolorbox{resultbox}[1]{
  enhanced,breakable,colback=DDTAGreen,colframe=DDTAGreenBorder,
  boxrule=0.8pt,arc=1.5mm,left=3mm,right=3mm,top=2.5mm,bottom=2.5mm,
  title={#1},fonttitle=\bfseries\sffamily,fontupper=\small\sffamily
}
\newtcolorbox{openbox}[1]{
  enhanced,breakable,colback=DDTAWarn,colframe=DDTAWarnBorder,
  boxrule=0.8pt,arc=1.5mm,left=3mm,right=3mm,top=2.5mm,bottom=2.5mm,
  title={#1},fonttitle=\bfseries\sffamily,fontupper=\small\sffamily
}

\lstset{
  basicstyle=\ttfamily\small,
  frame=single,
  framerule=0.4pt,
  rulecolor=\color{DDTABorder},
  backgroundcolor=\color{DDTAGray},
  breaklines=true,
  columns=fullflexible,
  keepspaces=true,
  showstringspaces=false,
  aboveskip=0.8em,
  belowskip=0.8em
}

\newcommand{\closed}{\textbf{CLOSED}}
\newcommand{\candidate}{\textbf{CANDIDATE}}
\newcommand{\strongcandidate}{\textbf{STRONG CANDIDATE}}
\newcommand{\openstatus}{\textbf{OPEN}}
\newcommand{\deferred}{\textbf{DEFERRED}}
\newcommand{\artifact}[1]{\nolinkurl{#1}}

\begin{document}
\pagenumbering{gobble}

\begin{titlepage}
\centering
\vspace*{2.1cm}
{\Large Documentation-Driven Threat Analysis\par}
\vspace{0.8cm}
{\Huge\bfseries S1.5 Working Note\par}
\vspace{0.35cm}
{\LARGE Requirement abstraction closure and provenance constraints\par}
\vspace{1.0cm}
{\large Formalizzazione temporanea tra SpecializedRequirement S1 e SecurityRequirement S2\par}
\vfill
\begin{minipage}{0.84\textwidth}
\small
\textbf{Stato.} Materiale di ricerca temporaneo. Non costituisce un capitolo autonomo della tesi, non promuove ancora una nuova foundation canonica e non chiude il modello di AnalysisRecord, Finding, revision o change event.

\medskip
\textbf{Baseline repository for this correction.} \artifact{nballestriero/documentation-driven-threat-analysis}, commit \artifact{544ef14d92b50cf956cf6e8dd1000079d757c49c}.

\medskip
\textbf{Scopo.} Chiudere l'astrazione comune \texttt{Requirement [abstract]}, registrare il rigetto della metaclasse \texttt{NormativeObligation} e preservare i vincoli minimi che il futuro modello di provenance dovra' soddisfare.
\end{minipage}
\vfill
{\large 14 agosto 2026\par}
\end{titlepage}

\pagenumbering{roman}
\tableofcontents
\newpage
\pagenumbering{arabic}

\section{Perche' esiste questa nota}

S1 ha chiuso la semantica di \texttt{SpecializedRequirement} come rafforzamento normativo governato di esattamente un \texttt{FunctionalRequirement}. La baseline corrente mantiene tuttavia una asimmetria strutturale:

\begin{lstlisting}
FunctionalRequirement
    functionalObligation
    functionalClause 1..*

SpecializedRequirement
    normativeObligation
\end{lstlisting}

La prima domanda S1.5 e' quindi piu' stretta della semantica di SecurityRequirement:

\begin{quote}
\textbf{FunctionalRequirement e SpecializedRequirement condividono un contratto astratto comune di Requirement?}
\end{quote}

In parallelo e' emersa una seconda domanda, distinta:

\begin{quote}
\textbf{Come deve DDTA preservare il fatto che una analisi abbia contribuito all'introduzione, revisione, sostituzione o retirement di documentazione governata?}
\end{quote}

Le due domande non devono essere fuse:

\begin{lstlisting}
requirement semantics
        !=
analysis/change provenance
\end{lstlisting}

\section{Collocazione futura nella tesi}

Questo materiale non giustifica un capitolo autonomo. La consolidazione piu' plausibile e' una sezione del Capitolo 4 dedicata alla struttura comune della documentazione governata e, al suo interno o immediatamente dopo, una sottosezione sull'astrazione comune dei Requirement.

\begin{figure}[h]
\centering
\begin{tikzpicture}[
  node distance=8mm,
  box/.style={draw=DDTADark,rounded corners,align=center,inner sep=6pt,minimum width=72mm},
  sub/.style={draw=DDTADark,rounded corners,align=left,inner sep=6pt,minimum width=72mm},
  arr/.style={-{Latex[length=2.4mm]},thick}
]
\node[box] (ch) {Capitolo 4 - Documentation metamodel e authoring};
\node[sub,below=of ch] (gd) {Contratto comune dei GovernedDocument\\\small identity, lifecycle, history/provenance [partly OPEN]};
\node[sub,below=of gd] (rq) {Requirement abstraction\\\small Requirement [abstract], normativeClause 1..*};
\node[sub,below=of rq] (types) {MacroRequirement - Decision - FunctionalRequirement - SpecializedRequirement};
\draw[arr] (ch)--(gd);\draw[arr] (gd)--(rq);\draw[arr] (rq)--(types);
\end{tikzpicture}
\caption{Collocazione futura proposta. Requirement non coincide con GovernedDocument.}
\end{figure}

La distinzione da preservare e':
\begin{itemize}
  \item \texttt{GovernedDocument} possiede proprieta' trasversali a tutti i documenti governati;
  \item \texttt{Requirement} accomuna soltanto i documenti che stabiliscono una obbligazione normativa;
  \item history/provenance e' piu' generale di Requirement e non deve essere confinata a SpecializedRequirement.
\end{itemize}

\newpage
\part*{Parte A - Requirement abstraction}
\addcontentsline{toc}{section}{Parte A - Requirement abstraction}

\section{Ipotesi S1.5-A e domanda di closure}

La gerarchia sottoposta a falsificazione e':

\begin{figure}[h]
\centering
\begin{tikzpicture}[
  node distance=9mm and 8mm,
  box/.style={draw=DDTADark,rounded corners,align=center,inner sep=6pt,text width=42mm},
  arr/.style={-{Latex[length=2.5mm]},thick}
]
\node[box] (gd) {GovernedDocument};
\node[box,below=of gd] (req) {Requirement\\\textit{abstract}};
\node[box,below left=of req] (fr) {FunctionalRequirement};
\node[box,below right=of req] (sr) {SpecializedRequirement\\\textit{abstract}};
\draw[arr] (gd)--(req);\draw[arr] (req)--(fr);\draw[arr] (req)--(sr);
\end{tikzpicture}
\caption{Common abstraction closed by S1.5-A.}
\end{figure}

La proprieta' comune e' semantica, non una mera uniformita' editoriale:

\begin{quote}
\textbf{Un Requirement e' una obbligazione normativa governata il cui contenuto normativo costituisce esattamente una unita' di requisito coerente.}
\end{quote}

Forma minima:

\begin{lstlisting}
Requirement [abstract]
    extends GovernedDocument

    normativeClause : NormativeClause [1..*]
\end{lstlisting}

Il Requirement e' esso stesso l'obbligazione normativa governata. Non necessita di possedere un secondo oggetto semantico che ripeta la stessa responsabilita'.

\section{Alternativa rigettata - metaclasse NormativeObligation}

La precedente ipotesi S1.5 introduceva:

\begin{lstlisting}
Requirement
    normativeObligation : NormativeObligation [1]

NormativeObligation
    normativeClause : NormativeClause [1..*]
\end{lstlisting}

La metaclasse aggiuntiva non possiede, nel baseline corrente, identity, lifecycle, governance, parentage, comportamento o semantica evolutiva indipendenti. La sua sola funzione sarebbe racchiudere le clausole normative gia' possedute dal Requirement.

\begin{openbox}{REJECTED - NormativeObligation metaclass}
La forma \texttt{Requirement HAS-A NormativeObligation} e' rigettata dal baseline L1 corrente. Il termine ``normative obligation'' rimane linguaggio definitorio: \texttt{Requirement IS-A governed normative obligation}.
\end{openbox}

\texttt{NormativeClause} rappresenta una o piu' clausole necessarie a esprimere completamente il Requirement. La sua esatta rappresentazione L2 testuale/strutturata resta separata dall'identita' L1 del Requirement.

\section{Coherent-unit invariant}

Un Requirement puo' richiedere piu' di una clausola testuale, ma le clausole devono esprimere congiuntamente una sola obbligazione normativa coerente:

\begin{lstlisting}
1 Requirement != 1 textual sentence
1 Requirement  = 1 coherent normative obligation
\end{lstlisting}

\begin{resultbox}{CLOSED - Requirement coherent-unit rule}
Tutte le normative clause possedute da un Requirement MUST esprimere congiuntamente esattamente una obbligazione normativa coerente.
\end{resultbox}

\textbf{Split test.} Se una clausola puo' essere introdotta, revisionata, ritirata o valutata indipendentemente senza cambiare l'identita' normativa delle clausole rimanenti, deve essere valutata come Requirement separato anziche' come ulteriore clausola dello stesso Requirement.

Questo evita di aggregare proprieta' specializzate indipendenti soltanto perche' condividono lo stesso parent FR.

\section{FunctionalRequirement come specializzazione}

Forma chiusa:

\begin{lstlisting}
FunctionalRequirement
    extends Requirement

    parentDecision : Decision [1]
\end{lstlisting}

Le clausole normative ereditate restano soggette agli invarianti FR gia' chiusi:
\begin{itemize}
  \item operationality;
  \item independent assessability;
  \item un solo comportamento, servizio, transizione o risultato osservabile coerente;
  \item completezza della correttezza funzionale ordinaria;
  \item implementation independence;
  \item normative prose primary.
\end{itemize}

Il precedente label \texttt{functionalObligation} viene normalizzato a L1 nella semantica comune di Requirement. Una rappresentazione L2 puo' conservare label FR-specifiche per l'authoring, pur mantenendo una sola autorita' semantica.

\section{SpecializedRequirement come specializzazione}

Forma chiusa:

\begin{lstlisting}
SpecializedRequirement [abstract]
    extends Requirement

    parentFunctionalRequirement : FunctionalRequirement [1]
\end{lstlisting}

Le clausole normative ereditate restano soggette alle regole CLOSED S1:
\begin{itemize}
  \item normative strengthening di esattamente un FR;
  \item composizione congiuntiva con il parent FR;
  \item removal test con preservazione della funzione ordinaria;
  \item nessuna duplicazione della correttezza funzionale;
  \item concern-specific obligation;
  \item positive subordinate action ammessa;
  \item autonomous capability boundary verso FR;
  \item realization independence;
  \item parent dependence.
\end{itemize}

Il contratto Requirement non indebolisce ne' sostituisce la semantica S1.

\section{Regression probes e negative controls}

\small
\begin{longtable}{p{0.24\textwidth}p{0.49\textwidth}p{0.15\textwidth}}
\toprule
\textbf{Probe} & \textbf{Osservazione} & \textbf{Esito} \\
\midrule
\endfirsthead
\toprule
\textbf{Probe} & \textbf{Osservazione} & \textbf{Esito} \\
\midrule
\endhead
FR-3.4 Deliver RecognitionCapture & Resta un Requirement operativo autonomamente significativo; piu' clausole possono esprimere correlation, delivery e failure semantics come unita' coerente. & PASS \\
\addlinespace
SR-3.4-C Confidentiality & Rafforza FR-3.4 e puo' essere espresso direttamente mediante clause senza oggetto NormativeObligation intermedio. & PASS \\
\addlinespace
SR-3.4-I Integrity & Stessa osservazione: nessuna perdita semantica eliminando la metaclasse wrapper. & PASS \\
\addlinespace
SR-3.4-P Authorized provenance & Acceptance condition e failure outcome possono restare piu' clause dello stesso SR finche' inseparabili dalla stessa obbligazione di provenance. & PASS \\
\addlinespace
Confidentiality vs authorized provenance & Possono evolvere, essere soddisfatte o diventare non applicabili indipendentemente: restano SR separati, non clause di un unico security Requirement. & PASS split \\
\addlinespace
PerformanceRequirement illustrativo & Un target prestazionale concern-specific puo' usare lo stesso contratto di clause mantenendo semantica di subtype. & PASS generality \\
\addlinespace
RegulatoryRequirement illustrativo & Un obbligo normativo/compliance puo' usare lo stesso contratto comune senza diventare FR. & PASS generality \\
\bottomrule
\end{longtable}
\normalsize

\subsection{Negative control - Decision}

Una Decision e' governata e puo' contenere linguaggio prescrittivo, ma registra un commitment significativo che restringe un MR. Non e' una obbligazione di soddisfazione funzionale o specializzata.

\begin{lstlisting}
Decision IS-NOT-A Requirement
\end{lstlisting}

\begin{resultbox}{NEGATIVE CONTROL - PASS}
L'astrazione Requirement non si espande automaticamente a ogni GovernedDocument.
\end{resultbox}

\section{CLOSED - S1.5-A Requirement abstraction}

\begin{resultbox}{CLOSED - Requirement abstraction}
\texttt{Requirement [abstract]} e' la superclass comune di FunctionalRequirement e SpecializedRequirement. Il Requirement e' esso stesso l'obbligazione normativa governata ed e' espresso da una o piu' normative clause che devono formare una sola unita' normativa coerente.
\end{resultbox}

Forma chiusa:

\begin{lstlisting}
Requirement [abstract]
    extends GovernedDocument
    normativeClause : NormativeClause [1..*]

FunctionalRequirement
    extends Requirement
    parentDecision : Decision [1]

SpecializedRequirement [abstract]
    extends Requirement
    parentFunctionalRequirement : FunctionalRequirement [1]
\end{lstlisting}

Invarianti chiusi:
\begin{enumerate}
  \item \textbf{Requirement identity}: il Requirement e' l'obbligazione normativa governata; non e' richiesta una metaclasse \texttt{NormativeObligation} separata.
  \item \textbf{Clause multiplicity}: un Requirement possiede una o piu' normative clause.
  \item \textbf{Coherent unit}: tutte le clause formano esattamente una obbligazione normativa coerente.
  \item \textbf{Split on independence}: obblighi indipendentemente evolvibili/assessable sono Requirement distinti.
  \item \textbf{Subtype preservation}: FR e SR conservano tutti i propri invarianti specifici.
  \item \textbf{Non-requirement boundary}: Decision e altri documenti governati non diventano Requirement soltanto perche' contengono linguaggio normativo.
\end{enumerate}

\begin{openbox}{REJECTED - L1 metaclass}
\texttt{NormativeObligation [metaclass]} e' rigettata: non aggiunge una responsabilita' semantica indipendente nel modello corrente.
\end{openbox}

\textbf{Reopen criterion.} Riaprire S1.5-A soltanto se un concreto subtype futuro non puo' usare \texttt{normativeClause [1..*]} piu' coherent-unit invariant senza perdita/distorsione, oppure se emerge la necessita' indipendente di identity/lifecycle per una obbligazione separata dal Requirement.

\newpage
\part*{Parte B - Provenance constraints}
\addcontentsline{toc}{section}{Parte B - Provenance constraints}

\section{Confine ereditato da S1}

S1 ha gia' distinto:

\begin{lstlisting}
provenance = why an obligation/change exists
semantics  = what the governed obligation requires
\end{lstlisting}

Il meccanismo formale di provenance era deliberatamente deferred. Anche la change-impact semantics gia' anticipa esiti di review quali:

\begin{lstlisting}
retain
revise
replace
non-applicable
\end{lstlisting}

S1.5 deve quindi registrare il requisito futuro del modello, senza chiudere prematuramente AnalysisRecord, Finding, revision history o change-event machinery.

\section{Analysis-to-document provenance requirement}

\begin{statusbox}{CANDIDATE - provenance requirement}
DDTA deve preservare provenance governata sufficiente a determinare quali cambiamenti accettati della documentazione siano stati motivati da ciascuna esecuzione di analisi e, inversamente, quali esecuzioni di analisi abbiano contribuito all'introduzione, revisione, sostituzione o retirement di ciascun documento governato.
\end{statusbox}

Il requisito e' potenzialmente applicabile a:

\begin{lstlisting}
MacroRequirement
Decision
FunctionalRequirement
SpecializedRequirement
\end{lstlisting}

Il caso facial-access \texttt{MR-0004} mostra gia' perche' la provenance futura non puo' essere confinata ai soli Requirement: una review downstream ha motivato la revisione della decomposizione MR.

\section{Invarianti di provenance preservabili ora}

\subsection{P1 - Provenance separation}

\begin{resultbox}{CLOSED candidate P1}
La provenance che spiega perche' un Requirement o altro documento governato e' stato introdotto, rivisto, sostituito o ritirato non appartiene alla sua obbligazione normativa.
\end{resultbox}

Ne segue che un campo intrinseco obbligatorio del tipo:

\begin{lstlisting}
SecurityRequirement.analysisId : exactly 1
\end{lstlisting}

non e' sufficiente ne' corretto come semantica del Requirement.

\subsection{P2 - Historical multiplicity}

\begin{resultbox}{CLOSED candidate P2}
Il futuro modello di provenance deve supportare una singola analisi con effetti su zero o molti documenti governati e un singolo documento governato con cambiamenti motivati da zero o molte analisi nel tempo.
\end{resultbox}

\begin{lstlisting}
one analysis execution
    -> effects on zero or many governed documents

one governed document
    -> effects motivated by zero or many analyses over time
\end{lstlisting}

\subsection{P3 - Analysis is not governance authority}

\begin{resultbox}{CLOSED candidate P3}
Una analisi puo' motivare o proporre una modifica, ma non canonizza automaticamente documentazione governata.
\end{resultbox}

\begin{center}
\begin{tikzpicture}[
  node distance=9mm,
  box/.style={draw=DDTADark,rounded corners,align=center,inner sep=6pt,minimum width=58mm},
  arr/.style={-{Latex[length=2.5mm]},thick}
]
\node[box] (a) {analysis / review};
\node[box,below=of a] (g) {governed acceptance / change};
\node[box,below=of g] (d) {governed document};
\draw[arr] (a)--node[right,font=\small]{motivates / proposes}(g);
\draw[arr] (g)--node[right,font=\small]{changes}(d);
\end{tikzpicture}
\end{center}

\subsection{P4 - Analyzed is not changed}

\begin{resultbox}{CLOSED candidate P4}
Il futuro modello deve distinguere un documento incluso nello scope di analisi da un documento effettivamente modificato, dopo governance review, a causa dell'analisi.
\end{resultbox}

Una analisi che esamina FR-X e non produce alcuna modifica accettata non deve creare una falsa storia di cambiamento.

\section{Meccanismo strutturale deliberatamente OPEN}

La seguente forma e' soltanto illustrativa:

\begin{lstlisting}
AnalysisRun
      |
      v
GovernedDocumentChangeEvent
      |
      v
GovernedDocument
\end{lstlisting}

Possibili event kind:

\begin{lstlisting}
INTRODUCED
REVISED
REPLACED
RETIRED
\end{lstlisting}

\begin{openbox}{OPEN - structural mechanism}
S1.5 non chiude AnalysisRun, AnalysisRecord, Finding, GovernedDocumentChangeEvent, RequirementRevision, event ordering, supersession semantics, analysis scope/coverage, governance acceptance workflow o le cardinalita' esatte fra risultati analitici e cambiamenti accettati.
\end{openbox}

Questi concetti dovranno essere falsificati insieme al futuro modello di analisi.

\section{Acceptance queries per il futuro modello}

Qualunque struttura venga scelta dovra' permettere entrambe le direzioni di interrogazione.

\subsection{Analysis verso accepted changes}

\begin{lstlisting}
AN-0042
    +-- introduced SR-10
    +-- introduced SR-11
    +-- contributed to revision of FR-7
    `-- contributed to retirement of SR-3
\end{lstlisting}

\subsection{Governed document verso analysis lineage}

\begin{lstlisting}
SR-10
    +-- introduced after AN-0042
    +-- revised after AN-0071
    `-- retired after AN-0108
\end{lstlisting}

Il modello deve inoltre consentire a una analisi di evidenziare un FunctionalRequirement ordinario mancante, senza costringere ogni output analitico a diventare SpecializedRequirement.

\section{Stato epistemico S1.5 dopo la correzione}

\small
\begin{longtable}{p{0.21\textwidth}p{0.69\textwidth}}
\toprule
\textbf{Stato} & \textbf{Elementi} \\
\midrule
\endfirsthead
\toprule
\textbf{Stato} & \textbf{Elementi} \\
\midrule
\endhead
\closed & Semantica operativa FR; semantica S1 SpecializedRequirement; SR esattamente un FR parent; provenance distinta dalla normative semantics; analisi non autorita' automatica della documentazione; \texttt{Requirement [abstract]} come superclass comune FR/SR; \texttt{Requirement.normativeClause [1..*]}; coherent-unit e split-on-independence invariants. \\
\addlinespace
\textbf{REJECTED} & \texttt{NormativeObligation} come metaclasse L1 separata; il termine resta linguaggio definitorio della semantica Requirement. \\
\addlinespace
\candidate & Historical many-to-many analysis/change lineage; provenance piu' ampia di Requirement; analyzed $\neq$ changed; analysis motivates/proposes, governance accepts. \\
\addlinespace
\openstatus & Contratto completo GovernedDocument; exact lifecycle/history model; ChangeEvent; RequirementRevision; AnalysisRun/AnalysisRecord; analysis type; change-kind vocabulary; supersession/retirement; structured applicability; rappresentazione L2 esatta delle normative clause. \\
\addlinespace
\deferred & Finding; Base Analysis/BAE; STRIDE; SecurityRequirement-specific provenance fields; controls; verification evidence. \\
\bottomrule
\end{longtable}
\normalsize

\section{Prossimo microstep raccomandato}

S1.5-A e' chiuso per il baseline corrente. I vincoli di provenance sono conservati, mentre il loro meccanismo strutturale resta OPEN.

Il prossimo microstep di ricerca e' quindi S2:

\begin{lstlisting}
SecurityRequirement IS-A SpecializedRequirement
\end{lstlisting}

S2 deve determinare soltanto la semantica e l'eventuale contratto minimo aggiuntivo che rende uno SpecializedRequirement specificamente un SecurityRequirement. Finding, Base Analysis, STRIDE, controls e provenance event model restano fuori dal microstep salvo che un controesempio concreto di SecurityRequirement costringa a correggere prima un layer proprietario.

\begin{statusbox}{Conservation note}
Questa nota rimane materiale di lavoro perche' la parte provenance e' ancora aperta. La sua futura consolidazione nel Capitolo 4 dovra' preservare la distinzione fra contratto comune dei GovernedDocument, astrazione Requirement chiusa in S1.5-A e provenance analitica ancora da formalizzare.
\end{statusbox}

\end{document}
